\documentclass[10pt,a4paper]{amsart}
\usepackage[margin=19mm]{geometry}
\usepackage{amsmath,amssymb,mathtools}
\usepackage[scr=rsfs,cal=euler]{mathalfa}
\usepackage{xfrac}
\usepackage[unicode,hidelinks,bookmarks=false]{hyperref}
\newcommand{\ie}{i.e.\ }
\newcommand{\cf}{cf.\ }
\newcommand{\resp}{respectively\ }
\newcommand{\Subsection}{\subsection}
\providecommand{\nobreakdash}{\nobreak\hbox{-}}
\setlength{\emergencystretch}{2em}
\multlinegap0pt
\newcommand{\RR}{\mathbb{R}}
\newcommand{\Sc}{\mathcal{S}}
\newcommand{\ZZ}{\mathbb{Z}}
\theoremstyle{plain}
\newtheorem{thm}{Theorem}[section]
\newtheorem{cor}[thm]{Corollary}
\newtheorem{lem}[thm]{Lemma}
\newtheorem{prop}[thm]{Proposition}
\newtheorem{prob}[thm]{Problem}
\newtheorem{q}[thm]{Question}
\newtheorem{conj}[thm]{Conjecture}
\newtheorem{assume}[thm]{Assumption}
\theoremstyle{definition}
\newtheorem{defi}[thm]{Definition}
\newtheorem{ex}[thm]{Example}
\theoremstyle{remark}
\newtheorem{rem}[thm]{Remark}
\title[Ehrhart polynomials of cyclic polytopes as averages of zonotope Ehrhart polynomials]{Ehrhart polynomials of cyclic polytopes as averages of zonotope Ehrhart polynomials}
\author{Masato Konoike}
\date{Source: arXiv:2609.28417v1 (CC BY 4.0)}
\begin{document}
\maketitle

\noindent\emph{Source and licence.} The delivered work is the article shown in the title above, version arXiv:2609.28417v1 (CC BY 4.0), distributed under the Creative Commons Attribution 4.0 International licence, as recorded in the arXiv licence metadata for this exact version and shown on the abstract page saved with this item. The full licence notice, which carries the licence URL and the address of that abstract page, is the bundled file \texttt{licenses/arxiv-2609.28417-CC-BY-4.0.txt}.
\par\smallskip

\noindent\textit{Editorial scope.} Counted unit: the article's Lemma (source label lem:zonotope-coefficients) (source0.tex lines 312--317, source label \texttt{lem:zonotope-coefficients}): ``For every , E Z d( s) (m) J n J d ( i J g i ) (s i:i J)m J .''. It is delivered together with the article's complete original proof (source0.tex lines 319--347, one proof environment adjacent to the statement, 29 lines), copied line for line from the article's own text. Required context retained uncounted: eq:stanley (lines 303--306). Editorial formatting only: an AMS \texttt{amsart} preamble that restates the article's own macro definitions and its own theorem declarations, and the theorem environments the retained lines use; the article's own class file is neither shipped nor input; the retained environments are renumbered sequentially in order of appearance, whereas the article prints them with its own numbering; the article's equation numbering is not reproduced and every cross-reference inside the retained text resolves to the wrapper's numbering. No citation occurs inside any retained line, so the wrapper carries no bibliography; All retained bodies are exact copies of the bundled \texttt{sources/211588/source0.tex}, so no omission receipt applies. No asset-inclusion command occurs inside any retained span and the package ships no image asset.


\section*{Retained context: eq:stanley (source0.tex lines 303--306)}

\begin{equation}\label{eq:stanley}
    E_{\sum_{i=1}^n[0,a_i]}(m)
    =\sum_{\substack{J\subseteq[n]\\(a_i)_{i\in J}\text{ independent}}} \mu(J)m^{|J|}.
\end{equation}


\section{The counted unit: Lemma (source label lem:zonotope-coefficients) and its complete original proof}

\begin{lem}\label{lem:zonotope-coefficients}
For every $\mathbf s\in\Sc$,
\begin{equation}\label{eq:zonotope-coefficients}
    E_{Z_d(\mathbf s)}(m)= \sum_{\substack{J\subseteq[n]\\|J|\le d}} \left(\prod_{i\in J}g_i\right) \Delta(s_i:i\in J)m^{|J|}.
\end{equation}
\end{lem}

\begin{proof}
Fix $J=\{i_1<\cdots<i_j\}\subseteq[n]$ with $1\le j\le d$. First consider the $d\times j$ matrix with columns
\[
u_d(s_{i_1}),\ldots,u_d(s_{i_j}).
\]
Its minor in the first $j$ rows is the Vandermonde determinant
\begin{equation}\label{eq:first-minor}
    \Delta(s_{i_1},\ldots,s_{i_j})>0.
\end{equation}

For any choice of $j$ rows, say $1\le r_1<\cdots<r_j\le d$, the corresponding minor is obtained by specializing
\[
f(X_1,\ldots,X_j)=\det\bigl((X_b^{r_a-1})_{1\le a,b\le j}\bigr)
\]
at $X_b = s_{i_b}$. This polynomial is alternating with integer coefficients. Therefore the Vandermonde polynomial
\[
\Delta(X_1,\ldots,X_j)=\prod_{a<b}(X_b-X_a)
\]
divides $f$ in $\ZZ[X_1,\ldots,X_j]$. After substituting $X_b=s_{i_b}$, every $j\times j$ minor is therefore an integer multiple of~\eqref{eq:first-minor}. Since~\eqref{eq:first-minor} itself is one of the minors, the greatest common divisor of all maximal minors is exactly
\[
\Delta(s_{i_1},\ldots,s_{i_j}).
\]

Restoring the factor $g_i$ in each column multiplies every maximal minor by $\prod_{i\in J}g_i$. Consequently,
\[
\mu(J)=\left(\prod_{i\in J}g_i\right)\Delta(s_i:i\in J).
\]
The nonzero minor \eqref{eq:first-minor} shows that every set of at most $d$ generators is linearly independent, while any set of more than $d$ generators is dependent since the generators lie in $\RR^d$. Together with $\mu(\emptyset)=\Delta(\emptyset)=1$, substituting the above multiplicities into \eqref{eq:stanley} gives \eqref{eq:zonotope-coefficients}.
\end{proof}

\end{document}
